\section{Grundlagen und verwandte Arbeiten}
\label{sec:related-work}

\subsection{Inhaltliche Stichpunkte}
\begin{itemize}
  \item \textbf{Architektur von Ladesäulen}
  \begin{itemize}
    \item Im Experiment repräsentieren Node A bis H acht Knoten.
    \item Die verwendeten Merkmale sind vier elektrische Messwerte und der
          Betriebszustand State (idle oder charging).
    \item Reale EVSE-Komponenten, Backend-Anbindung und Kommunikationspfade
          anhand geeigneter Quellen beschreiben; die Notebooks bilden keine
          vollständige physische EVSE-Architektur ab.
  \end{itemize}
  \item \textbf{Bedrohungsmodell}
  \begin{itemize}
    \item Zielklassen: Backdoor, syn-flood und ``none'' (benignes Verhalten).
    \item In Notebook~07 werden unter anderem Distraction Attack~1/2,
          Distributed SYN Flood, Worm Backdoor und Cluster Shuffling SYN Flood
          als Szenariogruppen ausgewertet.
  \end{itemize}
  \item \textbf{P2P-basierte Erkennung}
  \begin{itemize}
    \item Verwandte Ansätze zum dezentralen Informationsaustausch und zur
          verteilten Anomalieerkennung recherchieren und gegenüber dem
          hier ausgewerteten Verfahren abgrenzen.
    \item Konsensverfahren von einfacher gewichteter Merkmalsaggregation
          unterscheiden: Die Notebooks implementieren eine gewichtete
          Aggregation von Peer-Berichten, keinen dokumentierten
          Konsensalgorithmus.
  \end{itemize}
  \item \textbf{Datenerweiterung in der Cybersicherheit}
  \begin{itemize}
    \item TabGAN und Denoising Autoencoder (DAE) als Generatoren für
          strukturierte Daten einordnen.
    \item Motivation für synthetische Beispiele und Risiken hinsichtlich
          Realitätsnähe, Klassenabdeckung und Evaluationsverzerrung belegen.
  \end{itemize}
\end{itemize}

\subsection{Noch zu erklären}
\begin{itemize}
  \item Herkunft und Bedeutung der Labels Backdoor und syn-flood sowie die
        konkrete Durchführung der jeweiligen Angriffe.
  \item Warum elektrische Betriebsmerkmale und State diese Angriffe
        erkennbar machen sollen, obwohl die verwendeten Features keine
        Netzwerkpaketmerkmale sind.
  \item Welche EVSE-Standards, P2P-Ansätze und Verfahren zur synthetischen
        Datengenerierung für diese Arbeit relevant sind; passende Quellen
        und eine Abgrenzung zum eigenen Beitrag ergänzen.
\end{itemize}

\subsection{Grafikbedarf}
\begin{itemize}
  \item \textbf{Architekturskizze:} EVSE, lokale Sensorik/Erkennung und
        externe Kommunikationspartner; als konzeptionelle Darstellung
        kennzeichnen, falls nicht im Experiment implementiert.
  \item \textbf{Angriffsszenarien:} optional eine kleine Übersicht der
        Backdoor- und SYN-Flood-Szenarien aus Notebook~07.
\end{itemize}
